% !TeX spellcheck = en_US
%%%%%%%%%%%%%%%%%
% This is an example CV created using altacv.cls (v1.1.5, 1 December 2018) written by
% LianTze Lim (liantze@gmail.com), based on the
% Cv created by BusinessInsider at http://www.businessinsider.my/a-sample-resume-for-marissa-mayer-2016-7/?r=US&IR=T
%
%% It may be distributed and/or modified under the
%% conditions of the LaTeX Project Public License, either version 1.3
%% of this license or (at your option) any later version.
%% The latest version of this license is in
%%    http://www.latex-project.org/lppl.txt
%% and version 1.3 or later is part of all distributions of LaTeX
%% version 2003/12/01 or later.
%%%%%%%%%%%%%%%%

%% If you are using \orcid or academicons
%% icons, make sure you have the academicons
%% option here, and compile with XeLaTeX
%% or LuaLaTeX.
% \documentclass[10pt,a4paper,academicons]{altacv}

%% Use the "normalphoto" option if you want a normal photo instead of cropped to a circle
% \documentclass[10pt,a4paper,normalphoto]{altacv}

\documentclass[10pt,a4paper,ragged2e]{altacv}

%% AltaCV uses the fontawesome and academicon fonts
%% and packages.
%% See texdoc.net/pkg/fontawecome and http://texdoc.net/pkg/academicons for full list of symbols. You MUST compile with XeLaTeX or LuaLaTeX if you want to use academicons.

% Change the page layout if you need to
\geometry{left=1cm,right=9cm,marginparwidth=6.8cm,marginparsep=1.2cm,top=1.25cm,bottom=1.25cm}

% Change the font if you want to, depending on whether
% you're using pdflatex or xelatex/lualatex
\ifxetexorluatex
  % If using xelatex or lualatex:
  \setmainfont{Carlito}
\else
  % If using pdflatex:
  \usepackage[utf8]{inputenc}
  \usepackage[T1]{fontenc}
  \usepackage[default]{lato}
\fi

% Change the colours if you want to
\definecolor{VividPurple}{HTML}{084D6E}
\definecolor{SlateGrey}{HTML}{2E2E2E}
\definecolor{LightGrey}{HTML}{666666}
\colorlet{heading}{VividPurple}
\colorlet{accent}{VividPurple}
\colorlet{emphasis}{SlateGrey}
\colorlet{body}{LightGrey}

% Change the bullets for itemize and rating marker
% for \cvskill if you want to
\renewcommand{\itemmarker}{{\small\textbullet}}
\renewcommand{\ratingmarker}{\faCircle}

%% sample.bib contains your publications
%\addbibresource{sample.bib}

\begin{document}
\name{Gonçalo Pereira}
\tagline{Electrical and Computer Engineer}
% Cropped to square from https://en.wikipedia.org/wiki/Marissa_Mayer#/media/File:Marissa_Mayer_May_2014_(cropped).jpg, CC-BY 2.0
\photo{2.5cm}{me}
\personalinfo{%
  % Not all of these are required!
  % You can add your own with \printinfo{symbol}{detail}
  \email{goncalo\_pereira@outlook.pt}
  \phone{+351 919 157 890}
%  \mailaddress{Address, Street, 00000 County}
  \location{Porto, Portugal}
  \homepage{https://g-pereira.github.io/}
%  \twitter{@marissamayer}
  \linkedin{linkedin.com/in/g-pereira}
  \github{github.com/G-Pereira} % I'm just making this up though.
  %\skype{c4f4s0g0}
%   \orcid{orcid.org/0000-0000-0000-0000} % Obviously making this up too. If you want to use this field (and also other academicons symbols), add "academicons" option to \documentclass{altacv}
}

%% Make the header extend all the way to the right, if you want.
\begin{fullwidth}
\makecvheader
\end{fullwidth}

%% Depending on your tastes, you may want to make fonts of itemize environments slightly smaller
\AtBeginEnvironment{itemize}{\small}

%% Provide the file name containing the sidebar contents as an optional parameter to \cvsection.
%% You can always just use \marginpar{...} if you do
%% not need to align the top of the contents to any
%% \cvsection title in the "main" bar.
\cvsection[page1sidebar]{Experience}

\cvevent{Hardware Developer}{BOSCH Car Multimedia}{Jan 2023 -- Present}{Braga, Portugal}
\begin{itemize}
	\item Infotainment and Compute units hardware design for different vehicle manufacturers.
	\item Designer of several hardware modules integrated in several projects. 
	\item Responsibility for validation and certification processes (CE/FCC and OEM)
	\item Collaborative debug sessions with software team to solve issues -> Hardware/Software co-design
\end{itemize}

\divider

\cvevent{IoT Specialist}{BUILT CoLAB}{Mar 2021 -- Jan 2023}{Porto, Portugal}
\begin{itemize}
	\item Developed specialized sensor prototypes for construction machinery (e.g. trucks, pavers, rollers) to improve productivity of the construction sector by making it more digital to reduce overhead and improve observability of construction processes.
\end{itemize}

\divider

\cvevent{HDMI-LVDS Video Converter}{Summer Internship}{Jun 2019 -- Sep 2019}{APTIV, Braga, Portugal}
\begin{itemize}
	\item A video converter schematic circuit and layout were designed in order to be used as a debugging tool on \textbf{infotainment} systems development and validation.
	\item Several product life-cycle activities were performed, for example component replacement analysis due to components end-of-life as a way to interact with the company's team. 
\end{itemize}

\divider

\cvevent{Underwater Autonomous Vehicle for Target Searching}{Computer Vision}{Sep 2016 -- Dec 2017}{INESC TEC, Porto, Portugal}
\begin{itemize}
	\item I took part in a team which developed an underwater autonomous vehicle for a competition that had the mission to map and detect targets in an emergency response scenario where it's not proper to be explored by humans, inspired by the 2011 Fukushima accident.
\end{itemize}

\divider

\cvevent{IoT Sensorization Solution for Solar Panel Laboratory}{Field Experience}{Aug 2015 -- Jan 2016}{Lisbon, Portugal}
\begin{itemize}
	\item Developing a smart temperature and humidity sensor system for an oven to improve manufacturing yield was my initial contact with designing electronic systems to be used in a factory floor. 
\end{itemize}

\newpage

\cvsection[page2sidebar]{Activities}

\cvevent{LCD Porto - Laboratório de Criação Digital}{Coordinator and Member}{Nov 2016 -- Present}{Porto, Portugal}
\begin{itemize}
	\item Being involved in community and social technical projects has been giving me real world and a raw perspective where we as a society most need technology to be.
	\item Taught workshops about electronics and digital fabrication (ex. CAD, laser cutting, CNC)
	\item I develop and often provide support on electronics projects.
	\item As part of the coordination team, I have been developing a number of soft skill regarding management, organization and professional relationships. 
\end{itemize}

\divider

\cvevent{IEEE Robotics and Automation Student Chapter}{Member and Founder Chair}{Oct 2016 -- Sep 2019}{Porto, Portugal}
\begin{itemize}
	\item Being a leader in a team with people from different technical backgrounds improved my skills of knowledge sharing and organization of work teams by matching the right people together.
	\item Giving workshops also improved my experience in explaining topics clearly even if those were trivial for me.\\
	Most popular workshop: https://github.com/ieeeupsb/workshop\_ESP8266
\end{itemize}

\divider

\cvevent{IEEE R8 Student \& Young Professional Congress 2018, Porto}{Branding Team - Organization}{25 -- 29 Jul 2018}{Porto, Portugal}
\begin{itemize}
	\item My role was to arrange manufacture of marketing designs made by the team, balancing cost and quality of the different processes.
	\item During the event I also helped in logistics which gave me experience to anticipate needs of the participants and provide for them in proper time to make the event enjoyable and stress-free.
\end{itemize}

\divider

\cvevent{"Universidade Júnior" and "Mostra da Universidade do Porto"}{Tutor}{2016, 2017 and 2018}{FEUP, Porto, Portugal}
\begin{itemize}
	\item Talking about my study program and Engineering in general to help younger students choosing the right study path. To that end, workshops about robotics (sensors and actuators, maze solver and robot races) were presented.
\end{itemize}

\divider

\cvevent{"Projecto Saber"}{Volunteer}{Jan 2015 -- June 2015}{Vila Nova de Gaia, Portugal}
\begin{itemize}
	\item Helping younger students in their homework was my initial experience in knowledge sharing.
	\item Hours in record: 36
\end{itemize}

\newpage

\begin{fullwidth}

\cvsection{Publications}

\cvevent{Cooperative Perception for (Automated) Road Transport Enabled by Vehicular Networks}{2020 IEEE Vehicular Networking Conference (VNC)}{Dec 2020}{IT, Porto, Portugal}
\begin{itemize}
	\item A cooperative perception platform was designed based on COTS hardware, including an initial implementation of the ETSI TR 103 562 standard using ROS.
	\item Low level configurations were performed requiring to compile a Linux kernel with custom configurations (devicetree files).
	\item Different processing units and sensors were tested in order to understand the feasibility and latency of cooperative perception in real scenarios.
	\item This work was developed in the scope of my master thesis and the project \textit{POCI-01-0145-FEDER-016426}
\end{itemize}

\divider

\cvevent{Fuel Consumption Prediction for Construction Trucks: A Noninvasive Approach Using Dedicated Sensors and Machine Learning}{MDPI Infrastructures (Journal)}{Nov 2021}{BUILT CoLAB, Porto, Portugal}
\begin{itemize}
	\item A study of fuel consumption prediction based on route characteristics and vehicle movement data.
\end{itemize}

\end{fullwidth}

%\cvsection[page2sidebar]{a}

%\nocite{*}

%\printbibliography[heading=pubtype,title={\printinfo{\faBook}{Books}},type=book]

%\divider

%\printbibliography[heading=pubtype,title={\printinfo{\faFileTextO}{Journal Articles}}, type=article]

%\divider

%\printbibliography[heading=pubtype,title={\printinfo{\faGroup}{Conference Proceedings}},type=inproceedings]

%% If the NEXT page doesn't start with a \cvsection but you'd
%% still like to add a sidebar, then use this command on THIS
%% page to add it. The optional argument lets you pull up the
%% sidebar a bit so that it looks aligned with the top of the
%% main column.
%\addnextpagesidebar[-1ex]{page3sidebar}


\end{document}
